% \documentclass[landscape]{article}
\documentclass[]{article}
\usepackage{fontspec}
\usepackage[czech]{babel}
\setmainfont{TeX Gyre Heros}
\usepackage[margin=1.5cm]{geometry}
\usepackage{longtable}
\pagestyle{empty}

\usepackage{tikz}


% --------------------------------------------------
% Vodorovný regál
%
% #1 = číslo
% #2 = x
% #3 = y
% #4 = délka v gridových jednotkách
% --------------------------------------------------

\newcommand\regal[4]{%
  \draw[fill=gray!20]
    (#2,#3) rectangle ++(#4,1);
  \node [
    % circle,
    % fill=black,
    % text=white,
    % inner sep=2pt
    ]  at (#2+.5*#4,#3+.5) {\textbf{#1}};
}

% --------------------------------------------------
% Svislý regál
%
% #1 = číslo
% #2 = x
% #3 = y
% #4 = délka v gridových jednotkách
% --------------------------------------------------

\newcommand\regalv[4]{%
  \draw[fill=gray!20]
    (#2,#3) rectangle ++(1,#4);
    \node [
    % circle,
    % fill=black,
    % text=white,
    % inner sep=2pt
    ]  at (#2+.5,#3+.5*#4) {\textbf{#1}};
}

\newcommand\cisloregalu[1]{%
  \def\aktualniregal{#1}%
}

\makeatletter

\newcommand\nadpis[1]{%
  \gdef\aktualninadpis{#1}%
  \ifx\aktualniregal\@empty%
    \global\let\aktualniregal\predeslyregal%
  \fi%
}

\makeatother

\newcommand\radek[3]{%
  \aktualniregal &
  \aktualninadpis &
  #1 &
  #2
  \tabularnewline
  \let\predeslyregal\aktualniregal%
  \gdef\aktualniregal{}%
  \gdef\aktualninadpis{}%
}


\begin{document}

\parindent=0pt

\def\step{8mm}



\begin{center}
  \centering

{\Huge\bfseries Orientační plánek velké studovny}

\vfill

\begin{tikzpicture}[
  x=\step,
  y=\step
]

  \draw[step=\step,gray!15] (0,0) grid (12,26);
  \regal{1}{8}{25}{4}
  \regal{2}{8}{22}{4}
  \regal{3}{8}{21}{4}
  \regal{4}{8}{18}{4}
  \regal{5}{8}{17}{4}
  \regal{6}{8}{14}{4}
  \regal{7}{8}{13}{4}
  \regal{8}{8}{8}{4}
  \regal{9}{8}{7}{4}
  \regal{10}{8}{4}{4}
  \regal{11}{8}{3}{4}
  \regal{12}{8}{0}{4}
  \regalv{13}{0}{13}{3}
  \regalv{14}{0}{18}{3}

  % pomocný grid

  % vodorovné regály
  % \regal{1}{1}{7}{4}
  % \regal{2}{1}{5}{4}
  % \regal{3}{1}{3}{4}

  % další vodorovné regály
  % \regal{4}{7}{7}{4}
  % \regal{5}{7}{5}{4}
  % \regal{6}{7}{3}{4}

  % svislé regály
  % \regalv{7}{5}{1}{2}
  % \regalv{8}{7}{1}{2}

\end{tikzpicture}
\end{center}


\clearpage
\begin{longtable}{rllp{16em}}
\cisloregalu{1}
\nadpis{Biologie}
\radek{BIOL 1}{Obecná biologie}{}
\radek{BIOL 2}{Botanika}{}
\radek{BIOL 3}{Zoologie}{}
\radek{BIOL 4}{Environmentální studia}{}
\radek{BIOL 5}{Mikrobiologie a genetika}{}
\radek{BIOL 6}{Biologie člověka}{}
\radek{GEOL}{Geologie}{}
\radek{DID BIOL}{Didaktika biologie}{}
\radek{PŘÍR ZŠ}{Učebnice pro ZŠ Přírodověda}{}
\radek{BIOL SŠ}{Učebnice pro SŠ Biologie}{}

\cisloregalu{2}
\nadpis{Chemie}
\radek{CHEM}{Chemie}{}
\radek{DID CHEM}{Didaktika chemie}{}
\radek{CHEM ZŠ}{Učebnice pro ZŠ Chemie}{}
\radek{CHEM SŠ}{Učebnice pro SŠ Chemie}{}
\nadpis{Fyzika}
\radek{FYZ}{Fyzika}{}
\radek{FYZ ZŠ}{Učebnice pro ZŠ Fyzika}{}
\radek{FYZ SŠ}{Učebnice pro SŠ Fyzika}{}
\nadpis{Zeměpis}
\radek{ZEM}{Zeměpis}{}
\radek{ZEM ZŠ}{Učebnice pro ZŠ Zeměpis}{}
\radek{ZEM SŠ}{Učebnice pro SŠ Zeměpis}{}

\cisloregalu{3}
\nadpis{Matematika}
\radek{MAT}{Matematika}{}
\radek{DID MAT}{Didaktika matematiky}{}
\nadpis{Informatika}
\radek{IKM}{Informatika}{}

\cisloregalu{4}
\nadpis{Matematika}
\radek{MAT ZŠ}{Učebnice pro ZŠ Matematika}{}
\radek{MAT POM ZŠ}{Pomůcky pro ZŠ Matematika}{}
\radek{MAT SŠ}{Učebnice pro SŠ Matematika}{}
\nadpis{Informatika}
\radek{IKM ZŠ}{Učebnice pro ZŠ Informatika, IT knihy pro děti}{}

\cisloregalu{5}
\nadpis{Výtvarná výchova}
\radek{VV 1}{Umění a design}{}
\radek{VV 2}{České umění}{}
\radek{VV 3}{Výtvarné techniky}{}
\radek{VV 4}{Umění spolupráce}{}
\radek{VV 5}{Světové umění}{}
\radek{VV ZŠ}{Učebnice pro ZŠ Výtvarná výchova}{}
\radek{VV SŠ}{Učebnice pro SŠ Výtvarná výchova}{}

\cisloregalu{6}
\nadpis{Výtvarná výchova}
\radek{VV 6}{Architektura a urbanismus}{}
\radek{VV 7}{Dějiny umění}{}
\radek{DID VV}{Didaktika výtvarné výchovy}{}

\cisloregalu{7}
\nadpis{Hudební výchova}
\radek{HV 1}{Hudba}{}
\radek{HV 2}{Zpěv a sbory}{}
\radek{HV 3}{Hudební nástroje}{}
\radek{HV 4}{Hudební teorie}{}
\radek{DID HV}{Didaktika hudební výchovy}{}
\radek{HV ZŠ}{Učebnice pro ZŠ Hudební výchova}{}
\radek{HV SŠ}{Učebnice pro SŠ Hudební výchova}{}

\cisloregalu{8}
\nadpis{Tělesná výchova}
\radek{TV}{Tělesná výchova a sport}{}
\radek{DID TV}{Didaktika tělesné výchovy}{}
\nadpis{Výchova ke zdraví}
\radek{ZDRAV}{Výchova ke zdraví}{}

\cisloregalu{9}
\nadpis{Společenské vědy}
\radek{EKON}{Ekonomie}{}
\radek{POLIT}{Politologie}{}
\radek{NÁB}{Náboženství}{}
\radek{ZSV}{Základy společenských věd}{}
\radek{DID ZSV}{Didaktika základů společenských věd}{}
\nadpis{Informace a média}
\radek{INFO}{Informace, média}{}

\cisloregalu{10}
\nadpis{ZSV, občanská výchova}
\radek{OV ZŠ}{Učebnice pro ZŠ Člověk a jeho svět}{}
\radek{OV SŠ}{Učebnice pro SŠ Občanská výchova}{}

\cisloregalu{11}
\nadpis{Filozofie}
\radek{FIL 1}{Úvod do filozofie, etika}{}
\radek{FIL 2}{Dějiny filozofie}{}
\radek{FIL 3}{Česká filozofie}{}
\radek{FIL 4}{Světová filozofie}{}

\cisloregalu{12}
\nadpis{Sociologie}
\radek{SOCIOL 1}{Obecná sociologie}{}
\radek{SOCIOL 2}{Sociologické teorie}{}
\radek{SOCIOL 3}{Empirický výzkum – sociologie}{}
\radek{SOCIOL 4}{Sociologie vzdělávání}{}
\nadpis{Antropologie}
\radek{ANTROP}{Antropologie}{}
\nadpis{Psaní odborných textů}
\radek{ODB 1}{Psaní odborných textů}{}
\radek{ODB 2}{Metodologie výzkumu}{}

\cisloregalu{13}
\nadpis{Francouzština}
\radek{FR 1}{Jazykověda}{}
\radek{FR 2}{Literatura a literární věda}{}
\radek{FR 3}{Reálie}{}
\radek{FR 4}{Slovníky}{}
\radek{FR 5}{Příprava na zkoušky}{}
\radek{DID FR}{Didaktika francouzštiny}{}
\radek{FR ZŠ}{Učebnice pro ZŠ Francouzský jazyk}{}
\radek{FR SŠ}{Učebnice pro SŠ Francouzský jazyk}{}

\cisloregalu{14}
\nadpis{Latina}
\radek{LAT}{Latina}{}
\end{longtable}

\clearpage

\begin{center}
  \centering

{\Huge\bfseries Orientační plánek malé studovny}

\vfill

\begin{tikzpicture}[
  x=\step,
  y=\step,
  yscale=-1
]

  \draw[step=\step,gray!15] (0,0) grid (12,18);
  \regal{3}{1}{0}{11}
  \regalv{2}{0}{2}{3}
  \regal{1}{0}{6}{4}
  \regal{7}{0}{7}{4}
  \regal{4}{8}{6}{4}
  \regal{6}{8}{7}{4}
  \regalv{5}{5}{2}{5}
  \regalv{5}{6}{2}{5}
  \regal{9}{2}{10}{4}
  \regal{10}{2}{11}{4}
  \regal{8}{8}{10}{4}
  \regal{11}{8}{11}{4}
  \regal{12}{5}{14}{7}
  \regal{13}{5}{15}{7}
  \regal{15}{0}{18}{6}
  \regal{14}{7}{18}{5}

\end{tikzpicture}
\end{center}

\clearpage
\begin{longtable}{rllp{16em}}
\cisloregalu{1}
\nadpis{Speciální pedagogika}
\radek{SPPG 1}{Speciální pedagogika a inkluze}{}
\radek{SPPG 2}{Psychopedie (mentální postižení)}{}
\radek{SPPG 3}{Tyflopedie (zrakové postižení)}{}
\radek{SPPG 4}{Surdopedie (sluchové postižení)}{}
\radek{SPPG 5}{Specifické poruchy učení}{}
\cisloregalu{2}
\nadpis{Speciální pedagogika}
\radek{SPPG 6}{Narušená komunikační schopnost, logopedie}{}
\radek{SPPG 7}{Somatopedie (tělesné postižení)}{}
\radek{SPPG 8}{Etopedie (poruchy chování)}{}
\radek{SPPG 9}{Neurodivergence (autismus, ADHD)}{}
\radek{SPPG 10}{Sociální pedagogika (podpora, intervence)}{}
\radek{SPPG 11}{Kombinovaná postižení}{}
\radek{SPEC Š}{Učebnice a pomůcky pro speciální školy}{}
\cisloregalu{3}
\nadpis{Pedagogika}
\radek{PED 1}{Obecná pedagogika}{}
\radek{PED 2}{Obecná didaktika}{}
\radek{PED 3}{Pedagogické slovníky a encyklopedie}{}
\radek{PED 4}{Pedagogičtí pracovníci}{}
\radek{PED 5}{Komunikace ve škole, komunikace}{}
\radek{PED 6}{Předškolní a mimoškolní pedagogika}{}
\radek{MŠ POM}{MŠ – Pomůcky}{}
\radek{PED 7}{Vzdělávání dospělých}{}
\radek{PED 8}{Řízení a organizace školství, management}{}
\radek{PED 9}{Srovnávací pedagogika}{}
\radek{PED 10}{Dějiny pedagogiky}{}
\radek{PED 11}{Komeniologie}{}
\cisloregalu{4}
\nadpis{Psychologie}
\radek{PSYCH 1}{Obecná psychologie}{}
\radek{PSYCH 2}{Vývojová psychologie}{}
\radek{PSYCH 3}{Dějiny psychologie}{}
\radek{PSYCH 4}{Sociální psychologie}{}
\radek{PSYCH 5}{Psychologie osobnosti}{}
\radek{PSYCH 6}{Psychopatologie, klinická psychologie}{}
\radek{PSYCH 7}{Poradenství, psychoterapie}{}
\cisloregalu{5}
\nadpis{Psychologie}
\radek{PSYCH 7}{Poradenství, psychoterapie}{}
\radek{PSYCH 8}{Pedagogická psychologie}{}
\radek{PSYCH 9}{Psychologie – Ostatní}{}
\radek{PSYCH 10}{Řada Klasici psychologie}{}
\nadpis{Osobní rozvoj}
\radek{ROZ}{Osobní rozvoj}{}
\cisloregalu{6}
\nadpis{Český jazyk}
\radek{ČJ ZŠ}{Učebnice pro ZŠ Český jazyk}{}
\radek{ČJ POM ZŠ}{Pomůcky pro ZŠ Český jazyk}{}
\radek{ČJ SŠ}{Učebnice pro SŠ Český jazyk}{}
\cisloregalu{7}
\nadpis{Literatura}
\radek{LIT ZŠ}{Učebnice pro ZŠ Literární výchova}{}
\radek{ČÍT ZŠ}{Slabikáře, čítanky pro ZŠ}{}
\radek{LIT SŠ}{Učebnice pro SŠ Literatura}{}
\cisloregalu{8}
\nadpis{Český jazyk}
\radek{ČJ 1}{Mluvnice}{}
\radek{ČJ 2}{Pravopis}{}
\radek{ČJ 3}{Jazykověda}{}
\radek{ČJ 4}{Stylistika}{}
\radek{ČJ 5}{Slovníky a příručky}{}
\radek{DID ČJ}{Didaktika českého jazyka}{}
\cisloregalu{9}
\nadpis{Literatura}
\radek{LIT 1}{Teorie literatury}{}
\radek{LIT 2}{Čtenářství}{}
\radek{LIT 3}{Dějiny literatury}{}
\radek{LIT 4}{Literární encyklopedie, slovníky}{}
\radek{LIT 5}{Teorie literatury pro děti a mládež}{}
\radek{DID LIT}{Didaktika literatury}{}
\nadpis{Dramatická výchova}
\radek{DV 1}{Dramatická výchova}{}
\radek{DV 2}{Divadlo, film}{}
\cisloregalu{10}
\nadpis{Dějepis}
\radek{DĚJ ZŠ}{Učebnice pro ZŠ Dějepis}{}
\radek{DĚJ SŠ}{Učebnice pro SŠ Dějepis}{}
\radek{DĚJ 1}{Světové dějiny}{}
\cisloregalu{11}
\nadpis{Dějepis}
\radek{DĚJ 1}{Světové dějiny}{}
\radek{DĚJ 2}{České dějiny}{}
\radek{DĚJ 3}{Encyklopedie dějin}{}
\radek{DĚJ 4}{Dějiny UK a PedF}{}
\radek{DĚJ 5}{Pomocné vědy historické}{}
\radek{DID DĚJ}{Didaktika dějepisu}{}
\cisloregalu{12}
\nadpis{Beletrie}
\radek{BEL 1}{Próza}{}
\cisloregalu{13}
\nadpis{Beletrie}
\radek{BEL 1}{Próza}{}
\radek{BEL 2}{Poezie}{}
\radek{BEL 3}{Drama, libreta}{}
\cisloregalu{14}
\nadpis{Komiksy}
\radek{BEL 4}{Komiksy}{}
\nadpis{Literatura pro děti a mládež}
\radek{DĚT LIT 1}{Próza pro děti a mládež}{}
\cisloregalu{15}
\nadpis{Literatura pro děti a mládež}
\radek{DĚT LIT 1}{Próza pro děti a mládež}{}
\radek{DĚT LIT 2}{Poezie pro děti a mládež}{}
\radek{DĚT LIT 3}{Naučná literatura pro děti a mládež}{}
\end{longtable}





\end{document}
